\section{Materials and Methods}

\subsection{Cell Culture and Experimental Design}

Two cell lines were used: iMC ISOR544C (Line 1, isogenic control) and iMC MUTR544C (Line 2, mutant).

A full factorial design was implemented (Table~\ref{tab:experimental-design}), comprising 2 cell lines $\times$ 3 treatment conditions $\times$ 2 independent experiments, yielding 12 unique experimental conditions. Two independent experiments were conducted as biological replicates: EXP1 ($n = 158$ wells) and EXP2 ($n = 167$ wells), resulting in a total of 325 wound regions. Each well contained a single wound region, yielding one independent trajectory per well


\begin{table}[htbp]
	\centering
	\caption{Factorial experimental design.}
	\label{tab:experimental-design}
	\begin{tabular}{lcl}
		\toprule
		Factor & Levels & Description \\
		\midrule
		Cell line  & 2 & iMC ISOR544C, iMC MUTR544C \\
		Treatment  & 3 & DMSO, ALK5, Media \\
		Experiment & 2 & EXP1, EXP2 \\
		\bottomrule
	\end{tabular}
\end{table}

The ALK5 inhibitor was applied at 0.1~mM as a selective antagonist of the TGF-$\beta$ receptor type I (ALK5). Inhibition of TGF-$\beta$ signaling is relevant in this context as it promotes myogenic-to-fibrotic transdifferentiation \citep{Li2004}, suppresses myoblast differentiation \citep{McFarland2001}, and regulates leader cell dynamics during collective migration \citep{Vishwakarma2014}. DMSO was used as the vehicle control, while untreated media served as the basal reference condition.

Confluent monolayers were wounded using a standardized mechanical scratch assay. Treatment media were applied immediately after wounding. Time-lapse imaging began at $t = 0$ (baseline, fully open wound) and continued for approximately 24 hours, with images acquired at ~1-hour intervals, yielding 20–25 timepoints per trajectory. Images were captured at a resolution of approximately $900 \times 735$ pixels per field of view.

Images were acquired using a [microscope type] with a [X× objective] under [conditions].

\subsection{Automated Wound Segmentation: The Quantification Pipeline}

The Quantification pipeline operates in two stages: (1) a single-frame detector that extracts wound boundaries from each image independently, and (2) a temporal integrator that enforces the monotonic closure constraint across the time series. Parameters were selected based on empirical tuning on a held-out subset and fixed for all analyses.

For each trajectory, frames are processed sequentially as follows: each frame is independently segmented using the single-frame detector (Section~\ref{sec_single_frame}), producing validated edge estimates; these are refined by the temporal integrator (Section~\ref{sec:temporal-constraint}) under a monotonic constraint; the resulting edges are converted to wound masks and areas.

\subsubsection{Problem Formulation}

Let $I(x, y, t) \in [0, 1]$ denote the intensity-normalized grayscale image at spatial
position $(x, y)$ and discrete timepoint $t \in \{0, 1, \ldots, T\}$. The image has dimensions $H \times W$ (rows $\times$ columns).
We seek a binary wound mask:
\begin{equation}
  M(x, y, t) = \begin{cases}
    1 & \text{if pixel $(x, y)$ is wound at time $t$} \\
    0 & \text{otherwise}
  \end{cases}
  \label{eq:mask}
\end{equation}

and per-column edge functions:
\begin{align}
  u(x, t) &= \min\{y : M(x, y, t) = 1\} \quad \text{(upper edge)} \label{eq:upper-edge} \\
  \ell(x, t) &= \max\{y : M(x, y, t) = 1\} \quad \text{(lower edge)} \label{eq:lower-edge}
\end{align}

The wound area at time $t$ is:
\begin{equation}
  A(t) = \sum_{x=0}^{W-1} \max\!\bigl(0,\; \ell(x,t) - u(x,t) + 1\bigr)
  \label{eq:area}
\end{equation}

The biological axiom (monotonic closure) for wound area is a monotonically non-increasing
function of time:
\begin{equation}
  \boxed{A(t+1) \leq A(t) \quad \forall\; t \geq 0}
  \label{eq:monotonic}
\end{equation}

%We enforce this constraint at the per-column edge level (Section~\ref{sec:temporal-constraint}),
%which is strictly stronger than the aggregate area constraint.

\subsubsection{Single-Frame Detection}\label{sec_single_frame}

The single-frame detector maps an input image $I(x,y,t)$ to validated edge functions $u_\text{det}(x,t)$ and $\ell_\text{det}(x,t)$ through a seven-stage pipeline that transforms a raw grayscale image into a geometrically consistent wound mask with smooth boundary curves (Figure~\ref{fig:pipeline-example}).

\begin{figure}[H]
  \centering
  \includegraphics[width=0.9\linewidth]{methods_pipeline_example}
  \caption{%
    \textbf{Representative pipeline output on a single frame.}
    Phase-contrast image acquired at $t = 0$ (ALK5i condition, EXP2).
    The wound channel (center, low local variance) is flanked by cell monolayers
    (high local variance). Detected upper and lower boundaries are overlaid as
    colored lines after RANSAC outlier rejection and Savitzky--Golay smoothing.
    The shaded region corresponds to the reconstructed binary wound mask
    (Equation~\eqref{eq:mask-reconstruction}).%
  }
  \label{fig:pipeline-example}
\end{figure}

\textbf{Stage~1: Local Variance Texture Map.}
Cell monolayers exhibit high spatial variance due to cell borders, nuclei, and phase-contrast
halos, while the wound channel is comparatively smooth \citep{Yin2012} Figure~\ref{fig:pipeline-example}b).
We use local variance as a texture discriminator \citep{Koenig2017,Loizou2016,Ambrossio2022}.
For a square window of side length $w$ (default $w = 15$), the local variance at each pixel
is computed via the uniform-filter identity:
\begin{equation}
  \sigma^2(x, y) = \langle I^2 \rangle_w(x, y) - \bigl[\langle I \rangle_w(x, y)\bigr]^2
  \label{eq:local-variance}
\end{equation}
where $\langle \cdot \rangle_w$ denotes the box-filter mean over a $w \times w$ neighborhood,
yielding $O(HW)$ complexity independent of window size.

\textbf{Stage~2: Y-Band Detection.}
The wound is identified as a horizontal band spanning nearly the full image width.
We compute the row-averaged variance profile $\bar{V}(y)$, smooth it with a 1-D Gaussian
kernel ($\sigma_y = 10.0$ pixels), locate the global minimum $y_c$, and extract the
contiguous segment of sub-threshold rows (35th percentile) containing $y_c$,
expanded by a margin fraction of 0.15.

\textbf{Stage~3: Column-Wise Edge Extraction.}
Within the Y-band, upper and lower wound boundaries are extracted for each column.
The edge-detection threshold $\tau$ is derived adaptively from the cell regions flanking the
wound: $\tau = 0.5 \cdot \text{median}(\sigma^2_\text{cell})$.
Edge search is fully vectorized using array boolean operations.

\textbf{Stage~4: RANSAC Outlier Rejection.}
We apply an iterative polynomial-fit-and-reject scheme inspired by RANSAC
\citep{Fischler1981}.
A degree-2 polynomial captures genuine wound curvature while rejecting localized spikes.
Three iterations with a maximum deviation threshold of 25~pixels are used.

\textbf{Stage~5: Savitzky--Golay Smoothing.}
Residual high-frequency jitter is removed with a Savitzky--Golay filter \citep{Savitzky1964},
preserving edge shape without phase distortion.
Parameters: window length $w_\text{sg} = 51$, polynomial order $p_\text{sg} = 3$.

\textbf{Stage~6: Mask Reconstruction.}
The binary mask is reconstructed from smooth edges using vectorized broadcasting:
\begin{equation}
  M(x, y) = \bigl[u(x) \leq y \leq \ell(x)\bigr] 
  \label{eq:mask-reconstruction}
\end{equation}

\textbf{Stage~7: Geometric Validation.}
Four quality-control criteria are evaluated on each mask (Table~\ref{tab:qc-criteria}).

\begin{table}[htbp]
\centering
\caption{Geometric validation criteria applied to each wound mask.}
\label{tab:qc-criteria}
\begin{tabular}{llp{6.5cm}}
\toprule
Criterion & Threshold & Rationale \\
\midrule
Width coverage $r_w$       & $> 0.75$           & Wound must span most of image width \\
Aspect ratio $\rho$         & $> 1.5$            & Wound is wider than tall \\
Y-position $\bar{y}_\text{frac}$ & $\in (0.15, 0.85)$ & Wound is roughly centered \\
Thickness CV                & $< 0.60$           & Consistent thickness across columns \\
\bottomrule
\end{tabular}
\end{table}

A detection passes QC if and only if all four criteria are satisfied.
When the primary pipeline fails, an optional fallback strategy is activated using constrained
region growing from a seed at $y_c$ with morphological cleanup
(hole filling at 500~px, small-object removal at 1000~px, disk closing at radius~5).

\subsubsection{Temporal Constraint: Kalman-Filtered Monotonic Closure}
\label{sec:temporal-constraint}

A naive implementation of the monotonic closure constraint enforces a hard per-column clamp, such that the upper edge $u(x,t)$ can only move inward relative to $u(x,t-1)$ (and symmetrically for the lower edge). While this deterministic rule strictly encodes the biological axiom (Equation~\ref{eq:monotonic}), it introduces an irreversible failure mode: spurious inward detections—arising from debris, illumination artifacts, or transient focus shifts—are permanently propagated to all subsequent frames, with no mechanism for recovery.

To mitigate this limitation, we replace the hard clamp as the primary update mechanism with a per-column Kalman filter. The filter fuses the detector output with a monotonic prior, weighting each according to a data-driven estimate of observation uncertainty. Noisy or low-contrast detection are down-weighted, while consistent high-confidence observations are retained. 

A hard monotonic constraint is subsequently applied as a secondary safety mechanism to guarantee $A(t) \leq A(t-1)$ under all conditions. With this hierarchical design the Kalman filter handles typical artifact-contaminated observations, while the hard constraint provides an unconditional correctness for residual failure cases. 

The full per-frame processing sequence of the temporal integrator, including detector input, filtering steps, constraint enforcement, and numerical safeguards, is summarized in Supplementary Table~\ref{tab:temporal-sequence}. Four numerical safeguards (Supplementary Table~\ref{tab:temporal-sequence}, Steps 2, 7, 9, 11) prevent failure-mode propagation across frames; their formal statement follows below.

\textbf{Notation.}
Let $u_\text{det}(x, t)$ and $\ell_\text{det}(x, t)$ denote the raw single-frame detector
output, $u(x, t)$ and $\ell(x, t)$ the constrained edges,
and $P(x, t)$ the per-column Kalman uncertainty.

Let $u_\text{det}(x, t)$ and $\ell_\text{det}(x, t)$ denote the validated edge estimates produced by the single-frame detector (\ref{sec_single_frame}), following variance-based segmentation, outlier rejection, and geometric validation.


\textbf{Initialization} ($t = 0$):
\begin{equation}
  u(x, 0) = u_\text{det}(x, 0), \qquad \ell(x, 0) = \ell_\text{det}(x, 0)
\end{equation}
The $t = 0$ edges are accepted without constraint and serve as the global reference for each sequence.
The Kalman state is initialized to these values with initial uncertainty
$P(x, 0) = \max(R_\text{base},\; 1.0)$~px$^2$.

\textbf{Kalman-filtered constraint} ($t > 0$, primary mechanism).
The update at each frame proceeds in three sub-steps.

\emph{Prediction.}
The predicted state is the previous constrained edge; predicted uncertainty grows by process
noise $Q$:
\begin{equation}
  \hat{u}^-(x, t) = u(x, t-1), \qquad P^-(x, t) = P(x, t-1) + Q
  \label{eq:kalman-predict}
\end{equation}

Process noise $Q$ (px$^2$) is estimated empirically from QC-validated trajectories
$\mathcal{C}$ as the variance of consecutive frame-to-frame edge differences:
\begin{equation}
  \hat{Q} = \operatorname{Var}\!\Bigl(\bigl\{\,u_\text{det}(x,t{+}1) - u_\text{det}(x,t)\;\big|\; x,\,t \in \mathcal{C}\bigr\}\Bigr)
  \label{eq:Q-estimate}
\end{equation}
A conservative default $Q_0 = 400$~px$^2$ is used on first run.

\emph{Spatial pre-smoothing.}
The raw detected edge is smoothed with a 1-D Gaussian kernel ($\sigma_\text{spatial} = 5$
columns) before the Kalman update to suppress column-level artefacts:
\begin{equation}
  \tilde{u}_\text{det}(x, t) = \bigl(G_{\sigma} * u_\text{det}(\cdot,\, t)\bigr)(x)
  \label{eq:spatial-smooth}
\end{equation}

\emph{Per-column observation noise.}
Observation noise $R(x, t)$ is estimated from three data-driven signals:
\begin{equation}
  R(x, t) = R_\text{base}(t) \cdot C^{-2}(x, t) \cdot J(x, t) \cdot m_\text{method}
  \label{eq:R-estimate}
\end{equation}
where $R_\text{base}(t) = \operatorname{Var}_{x}(u_\text{det} - \tilde{u}_\text{det})$
is the frame-level detector noise, $C(x,t)$ is the local wound-vs-cell variance contrast
(low contrast implies high $R$), $J(x,t)$ is a jump penalty for large deviations from the
prediction, and $m_\text{method} \in \{1.0, 3.0\}$ penalises the less reliable fallback
detector.

\emph{Kalman update.}
The gain $K(x,t) \in [0, 1]$ blends prediction and observation:
\begin{align}
  K(x, t) &= \frac{P^-(x, t)}{P^-(x, t) + R(x, t)} \label{eq:kalman-gain} \\
  u_\text{filt}(x, t) &= u(x, t-1) + K(x, t)\bigl[\tilde{u}_\text{det}(x, t) - u(x, t-1)\bigr] \label{eq:kalman-update} \\
  P_\text{upd}(x, t) &= \bigl(1 - K(x, t)\bigr) \cdot P^-(x, t) \label{eq:P-update}
\end{align}
High $R$ (noisy observation) implies small $K$, causing the filter to rely on the monotonic
prior.
Low $R$ (clean, high-contrast frame) implies large $K$, trusting the detector.

\textbf{Hard safety net} (secondary guarantee).
After the Kalman blend, the biological invariant is enforced unconditionally:
\begin{align}
  u(x, t) &= \max\!\bigl(u_\text{filt}(x, t),\; u(x, t-1)\bigr) \label{eq:hard-upper} \\
  \ell(x, t) &= \min\!\bigl(\ell_\text{filt}(x, t),\; \ell(x, t-1)\bigr) \label{eq:hard-lower}
\end{align}
This guarantees $A(t) \leq A(t-1)$ for all $t > 0$ regardless of Kalman output.

\textbf{Wound closure freeze.}
When $\text{prev\_area} / A(0) < \theta_\text{close} = 0.01$, the wound is declared fully
closed and all subsequent frames inherit the preceding frame's edges.

\textbf{Numerical safeguards.}
Four safeguards prevent numerical failure from propagating across frames:
(i) non-finite observation noise $R$ is clamped to $Q_0 = 400$~px$^2$;
(ii) non-finite Kalman filter outputs are reset to the previous constrained edge;
(iii) per-column uncertainty $P$ is reset to $P^-$ when the safety net overrides the Kalman
estimate, preventing overconfidence propagation; and
(iv) a last-resort area comparison emits a \texttt{RuntimeWarning} and reverts to the previous
frame if $A(t) > A(t-1)$ persists after all constraint paths.

\textbf{Complexity.}
The Kalman update adds $O(W)$ operations per frame.
Total per-frame complexity remains $O(HW)$, dominated by the variance map computation.

\subsubsection{Trajectory Processing and Parallel Execution}

For each biological sample, images are sorted by acquisition time, fed sequentially to the
segmenter, and results serialized.
Independent samples are processed in parallel using
\texttt{concurrent.futures.ProcessPoolExecutor} with explicitly controlled worker count
($n_\text{workers}$, default 4, capped at available CPU cores).
Each worker instantiates its own segmenter with independent state, ensuring no shared mutable
data.

\subsection{Quality Control}

Quality control operates at three hierarchical levels:

\textbf{Level~1: Single-frame geometric validation} (Section~2.3.2, Stage~7).
Applied automatically to every frame; flags implausible geometry.

\textbf{Level~2: Baseline distributional analysis.}
The $t = 0$ mask undergoes additional scrutiny via row- and column-wise projection analysis,
requiring exactly one prominent peak and satisfying skewness and flatness criteria.

\textbf{Level~3: Manual expert curation.}
For high-stakes analyses, baseline images are reviewed by a human expert.
For a subset of images, the reviewer draws ground-truth polygon annotations enabling
pixel-level validation.

\subsection{Downstream Metrics}

\subsubsection{Fractional Wound Closure}

Fractional wound closure was defined as:
\begin{equation}
  f(t) = \frac{A(0) - A(t)}{A(0)}
  \label{eq:fractional-closure}
\end{equation}
ranging from 0 (no closure) to 1 (complete closure), where $A(t)$ is already strictly
non-increasing by construction.

\subsubsection{Edge Displacement and Closure Speed}

Per-edge displacements between $t = 0$ and target time $t_T$:
\begin{align}
  \Delta u(x) &= u(x, t_T) - u(x, 0) \label{eq:delta-upper} \\
  \Delta \ell(x) &= \ell(x, 0) - \ell(x, t_T) \label{eq:delta-lower}
\end{align}

Mean closure speed:
\begin{equation}
  v_\text{mean} = \frac{v_\text{upper} + v_\text{lower}}{2}, \quad
  v_\text{edge} = \frac{1}{|\mathcal{V}|} \sum_{x \in \mathcal{V}} \frac{|\Delta_\text{edge}(x)|}{t_T}
  \label{eq:closure-speed}
\end{equation}

\subsection{Statistical Analysis}

Pairwise treatment comparisons versus DMSO were performed using Welch's $t$-test
(no equal-variance assumption), stratified by cell line.
Multiple comparison correction was applied using the Bonferroni method.
Effect sizes were quantified as Cohen's $d$:
\begin{equation}
  d = \frac{\bar{x}_\text{treated} - \bar{x}_\text{control}}{s_\text{pooled}}
  \label{eq:cohens-d}
\end{equation}

The time-series model used a linear mixed-effects framework:
\begin{equation}
  \text{distance}_{ij} = \mu + \beta_1 t + \beta_2 \text{condition} + \beta_3 \text{cell\_line}
    + \beta_4 (t \times \text{condition}) + \beta_5 (t \times \text{cell\_line})
    + b_{0i} + b_{1i} t + \varepsilon_{ij}
  \label{eq:mixed-model}
\end{equation}
where $b_{0i}$ and $b_{1i}$ are random intercept and slope for trajectory $i$.
Normality of residuals was assessed by the Shapiro--Wilk test; homogeneity of variance was
assessed by Levene's test (see Supplementary Table~\ref{tab:supp-assumption-checks}).

\subsection{Implementation}

\begin{table}[htbp]
\centering
\caption{Software dependencies and versions.}
\label{tab:implementation}
\begin{tabular}{ll}
\toprule
Component & Version / Library \\
\midrule
Language        & Python 3.11 \\
Array operations & NumPy 1.24+ \\
Image I/O       & OpenCV 4.x \\
Filtering       & SciPy (\texttt{uniform\_filter}, \texttt{gaussian\_filter1d}, \texttt{savgol\_filter}) \\
Morphology      & scikit-image \citep{vanderWalt2014} \\
Data management & pandas, pickle \\
Visualization   & matplotlib \\
Parallelism     & \texttt{concurrent.futures.ProcessPoolExecutor} \\
\bottomrule
\end{tabular}
\end{table}

Software dependencies are listed in Table~\ref{tab:implementation}.
The pipeline is fully deterministic and reproducible.
Per-frame computational complexity is $O(HW)$, dominated by the variance map.
For a typical $1024 \times 1280$ image this corresponds to $\sim 1.3 \times 10^6$ operations
per stage.
